\section{Self-check}

\subsection{Parameters of a convolutional layer versus a fully connected layer}

\subsection{Why convolution and pooling are less sensitive to a small shift}

\subsection{Recognising overfitting from the training and validation curves}

\section{Training}

\begin{figure}[H]
  \centering
  \includegraphics[width=\textwidth]{stage3/curves.pdf}
  \caption{Training and validation loss (left) and accuracy (right) per epoch. The vertical line marks the epoch whose weights were kept.}
  \label{fig:s3-curves}
\end{figure}

\begin{table}[H]
  \centering
  \caption{Training summary of the CNN.}
  \label{tab:s3-training}
  \input{../tables/s3_training.tex}
\end{table}

\section{Comparison}

\begin{table}[H]
  \centering
  \caption{Test accuracy of the three models: softmax on raw pixels, softmax on PCA features, and the CNN.}
  \label{tab:s3-comparison}
  \input{../tables/s3_comparison.tex}
\end{table}

\section{Learned feature space}

\begin{figure}[H]
  \centering
  \begin{subfigure}[t]{0.48\textwidth}
    \includegraphics[width=\textwidth]{stage3/feature_space_raw.pdf}
    \caption{Raw pixels.}
  \end{subfigure}\hfill
  \begin{subfigure}[t]{0.48\textwidth}
    \includegraphics[width=\textwidth]{stage3/feature_space_cnn.pdf}
    \caption{CNN features before the softmax layer.}
  \end{subfigure}
  \caption{Validation images projected to two dimensions with PCA and coloured by letter.}
  \label{fig:s3-feature-space}
\end{figure}

\begin{table}[H]
  \centering
  \caption{Silhouette scores of the letter clusters in raw-pixel space and in the learned feature space.}
  \label{tab:s3-silhouette}
  \input{../tables/s3_silhouette.tex}
\end{table}

\section{Robustness}

\begin{table}[H]
  \centering
  \caption{Test accuracy of the three models on unshifted images and on images shifted by 1, 2 and 3 pixels (mean over the four directions).}
  \label{tab:s3-shift}
  \input{../tables/s3_shift.tex}
\end{table}

\begin{figure}[H]
  \centering
  \begin{subfigure}[t]{0.58\textwidth}
    \includegraphics[width=\textwidth]{stage3/shift_examples.pdf}
    \caption{A test image shifted by 0, 1, 2 and 3 pixels.}
  \end{subfigure}\hfill
  \begin{subfigure}[t]{0.38\textwidth}
    \includegraphics[width=\textwidth]{stage3/shift_acc.pdf}
    \caption{Test accuracy against shift size.}
  \end{subfigure}
  \caption{Effect of shifting the test images.}
  \label{fig:s3-shift}
\end{figure}

\section{Prediction vs.\ result}

\subsection{Question 1: do learned features do better?}

\subsection{Question 2: what did it learn?}

\subsection{Question 3: are learned features more robust?}

\section{AI-use note}

\appendix

\section{Shift results per direction}

\begin{table}[H]
  \centering
  \caption{Test accuracy for each shift direction separately.}
  \label{tab:s3-shift-directions}
  \input{../tables/s3_shift_directions.tex}
\end{table}
